% Adaptación del preámbulo de Paper III v1.5 español; no plantilla standalone de Pandoc.
\documentclass[spanish,11pt,a4paper]{article}
\usepackage[margin=1in]{geometry}
\usepackage{fontspec}
\setmainfont{lmroman10-regular.otf}[BoldFont=lmroman10-bold.otf,ItalicFont=lmroman10-italic.otf,BoldItalicFont=lmroman10-bolditalic.otf]
\setsansfont{lmsans10-regular.otf}[BoldFont=lmsans10-bold.otf]
\setmonofont{lmmono10-regular.otf}[Scale=MatchLowercase]
\newlength{\LeanCharWidth}
\usepackage{amsmath,amssymb}
\usepackage[spanish,es-noshorthands,es-tabla]{babel}
\usepackage{microtype,parskip}
\usepackage{graphicx,float,caption}
\usepackage{longtable,booktabs,array,calc}
\usepackage{xcolor}
\usepackage{xurl}
\usepackage{hyperref,bookmark}
\usepackage{etoolbox,needspace}
\setcounter{secnumdepth}{-\maxdimen}
\setlength{\emergencystretch}{3em}
\providecommand{\tightlist}{\setlength{\itemsep}{0pt}\setlength{\parskip}{0pt}}
\newcounter{none}
\AtBeginEnvironment{longtable}{\Needspace{8\baselineskip}\small}
\pretocmd{\section}{\Needspace{5\baselineskip}}{}{}
\pretocmd{\subsection}{\Needspace{4\baselineskip}}{}{}
\pretocmd{\subsubsection}{\Needspace{4\baselineskip}}{}{}
\captionsetup{font=small,labelfont=bf}
\urlstyle{same}
\hypersetup{colorlinks=true,linkcolor=blue,urlcolor=blue,citecolor=blue,pdflang={es-ES}}
\title{Estabilidad exacta de las particiones en cliques de grafos cordales casi
extremales}
\author{Juan Pablo Traverso Gianini\\
Investigador independiente, Santiago, Chile}
\date{}\begin{document}\raggedbottom\maketitle
Nota de investigación II de la serie, complementaria de los Papers IV y
VI\\
Preprint 1.0 --- 18 de octubre de 2026

\subsection{Resumen}\label{resumen}

Sea \(M(n)=\lfloor n(n+1)/6\rfloor\) y sea
\(d=M(n)-\operatorname{cp}(G)\) el déficit de un grafo cordal.
Combinamos el teorema estructural del Paper IV con la fórmula local
paramétrica del Paper VI. Para orden suficientemente grande y
\(0\le d\le n/2000\), los vértices de grado mayor que \(n/2\) forman un
clique \(P\). Respecto de este clique, el déficit es exactamente la suma
de la pérdida causada por su tamaño, el número de aristas cruzadas
ausentes y el doble del número de aristas exteriores. En particular, el
grafo difiere del grafo split completo asociado en a lo sumo \(d\)
aristas, y existe una partición óptima que utiliza cliques de orden a lo
sumo cuatro. La demostración obtiene primero una cota gruesa de edición
a partir de IV, la utiliza para entrar en el régimen local de VI y
después recupera la fórmula exacta. Se trata de una aplicación de los
dos papers, con un companion formal separado.

Una partición en cliques cubre cada arista exactamente una vez mediante
cliques de orden al menos dos. Su número mínimo de bloques se denota por
\(\operatorname{cp}(G)\). Escribimos \(\operatorname{cp}_{\le4}(G)\)
cuando los bloques se restringen a orden a lo sumo cuatro. Todos los
grafos son finitos y simples. Para conjuntos disjuntos de vértices
\(P,O\), la notación \(K_P\vee I_O\) denota el grafo split completo con
clique \(P\) y conjunto independiente \(O\).

Aquí importa distinguir una fórmula local de un teorema estructural. El
Paper VI da la igualdad exacta una vez que el grafo está localmente
cerca de un grafo split completo adecuado. El Paper IV proporciona ese
clique de referencia para un grafo cordal casi extremal. Esta nota
documenta cómo se combinan ambos enunciados y mantiene explícita la
identidad de déficit resultante. La cota eventual de Erdős 81 ya está
establecida en IV y en los otros desarrollos comparados en su Sección 8.
Esta nota trata la consecuencia de estabilidad enunciada, no una nueva
solución de ese problema.

\subsection{1. Estabilidad exacta en una ventana lineal de
déficit}\label{estabilidad-exacta-en-una-ventana-lineal-de-duxe9ficit}

Definamos

\[
M(n)=\left\lfloor\frac{n(n+1)}6\right\rfloor,
\qquad B(n,p)=p(n-p)-\binom p2,
\qquad \eta=10^{-16}.
\]

\Needspace{6\baselineskip}
\textbf{Teorema 1.1.} Utilizando los teoremas de redondeo uniforme en el
régimen lejano y de estructura cercana del Paper IV que se enuncian a
continuación, existe un entero \(N\) tal que todo grafo cordal \(G\) de
orden \(n\ge N\), con déficit

\[
d=M(n)-\operatorname{cp}(G),\qquad 0\le d\le n/2000,
\]

\Needspace{4\baselineskip}admite la siguiente descripción. El conjunto

\[
P=\{v:2d_G(v)>n\}
\]

es un clique. Escribamos \(p=|P|\) y \(O=V(G)\setminus P\), sea \(A\) el
número de aristas ausentes entre \(P\) y \(O\), y definamos
\(m=e(G[O])\). Entonces

\[
\boxed{\operatorname{cp}(G)=\operatorname{cp}_{\le4}(G)
=B(n,p)-A-2m,}
\tag{1}
\] \[
\boxed{d=\bigl(M(n)-B(n,p)\bigr)+A+2m,}
\tag{2}
\] \[
\boxed{|E(G)\triangle E(K_P\vee I_O)|=A+m\le d.}
\tag{3}
\]

Además,

\[
(6p-2n-1)^2\le24d+25.
\tag{4}
\]

La ecuación (2) separa el déficit en cantidades que pueden leerse del
grafo una vez conocido el clique. El primer término mide el efecto de su
tamaño; los otros cuentan las aristas cruzadas ausentes y las aristas
exteriores efectivas. La ecuación (3) acota entonces el número total de
ediciones respecto del grafo split completo asociado. El criterio de
grados identifica el mismo clique sin requerir una búsqueda entre
biparticiones de vértices.

El tamaño \(p\) no tiene que ser un tamaño maximizante, como
\(\lfloor(n+1)/3\rfloor\). Por ello se necesita la Proposición 2.1
paramétrica de VI, y no basta su Teorema A equilibrado.

\subsection{2. Resultados estructurales y de redondeo
utilizados}\label{resultados-estructurales-y-de-redondeo-utilizados}

Sea \(W^*(G)\) la ganancia fraccional máxima de un empaquetamiento de
triángulos y copias de \(K_4\), con ganancias 2 y 5. Escribamos
\(F_4(G)=e(G)-W^*(G)\).

\begin{enumerate}
\def\labelenumi{\arabic{enumi}.}
\tightlist
\item
  \textbf{Redondeo uniforme {[}1, Teorema 3.1{]}.} Para
  \(\zeta=\eta/2\), una vez que \(n\ge N(\eta/2)\), existe un
  empaquetamiento físico cuya completación en cliques tiene a lo sumo \[
  F_4(G)+(\eta/2)n^2
  \] bloques. Este resultado se aplica uniformemente a los grafos
  considerados, no solo a un optimizador fraccional previamente
  seleccionado.
\item
  \textbf{Testigo de estructura cercana {[}1, Proposición 4.3, Teorema
  5.0 y Lema 5.3{]}.} Si \(n\ge4\cdot10^{12}\), \(G\) es cordal y \[
  F_4(G)\ge n^2/6-\eta n^2,
  \] IV produce una raíz efectiva \(P\) que es un clique, un
  empaquetamiento físico y su partición completada de tamaño \(L\), con
  recuentos efectivos \(m,A\) de aristas exteriores y enlaces ausentes,
  que satisfacen \[
  L\le B(n,p)-m/20-A/2.
  \tag{5}
  \]
\end{enumerate}

En el primer resultado de entrada, completar un empaquetamiento
significa cubrir individualmente cada arista no utilizada. El objeto
resultante es una partición del grafo original, por lo que su número de
bloques es una cota superior de \(\operatorname{cp}(G)\). En el segundo
resultado, la raíz, el empaquetamiento, la partición y los recuentos
pertenecen a un único testigo. Utilizaremos su partición para controlar
los recuentos y su raíz para aplicar VI.

Los recuentos de este testigo son las aristas exteriores y las aristas
cruzadas ausentes efectivas de su raíz regularizada, que es un clique.
La ecuación (5.15) de IV da (5); su versión más fuerte (5.15a) se
refiere al mismo testigo. Esta identificación física es necesaria antes
de aplicar el teorema local de VI. El companion formal separado
documenta las interfaces fuente.

\subsection{3. El umbral de orden}\label{el-umbral-de-orden}

Elijamos \(N\) de modo que, para todo \(n\ge N\),

\[
n\ge N(\eta/2),\qquad n\ge4\cdot10^{12},\qquad n\ge100,
\qquad \frac\eta2n^2>\frac n{2000}+1.
\tag{6}
\]

Tal \(N\) existe porque \(\eta>0\). Para la última condición, puramente
numérica, \(n\ge2\cdot10^{13}\) ya basta; el umbral desconocido del
redondeo uniforme puede ser mayor. Es un teorema eventual, no una
afirmación numérica para todos los órdenes.

\subsection{4. Demostración}\label{demostraciuxf3n}

Primero excluimos el régimen lejano. Después utilizamos el testigo
cercano para acotar el número de aristas editadas y el tamaño de su raíz
clique. Estas dos cotas verifican las hipótesis de VI; solo después de
aplicar su fórmula exacta mejoramos la cota de edición e identificamos
la raíz por grados.

Supongamos primero que \(F_4(G)<n^2/6-\eta n^2\). El redondeo uniforme
da una partición física en cliques con menos de

\[
\frac{n^2}{6}-\frac\eta2n^2
<\frac{n^2}{6}-\frac n{2000}-1
\le M(n)-d=\operatorname{cp}(G)
\]

bloques, lo que es una contradicción. La última desigualdad no estricta
utiliza \(M(n)\ge n^2/6-1\) y \(d\le n/2000\). Por tanto, se aplica el
resultado de estructura cercana.

Utilicemos su raíz \(P\), sus recuentos efectivos \(m,A\) y su partición
de tamaño \(L\). La identidad cuadrática da
\(B(n,p)\le Q(n)=(2n+1)^2/24\). Para órdenes enteros, \(Q(n)-M(n)\) vale
\(1/24\) si \(n\equiv0,2\pmod3\), y \(9/24\) si \(n\equiv1\pmod3\). Por
tanto, \(Q(n)<M(n)+1\); como \(B(n,p)\) es entero, \(B(n,p)\le M(n)\).
Así,

\[
M(n)-d=\operatorname{cp}(G)\le L
\le B(n,p)-m/20-A/2\le M(n)-m/20-A/2.
\]

Las dos primeras conclusiones siguientes se obtienen comparando los
extremos de esta cadena. Para la tercera, multiplicamos la cota
\(m/20+A/2\le d\) por 20 y utilizamos \(A\ge0\):

\[
B(n,p)\ge M(n)-d,\qquad m/20+A/2\le d,
\qquad m+A\le20d.
\tag{7}
\]

Como esta raíz es un clique, el grafo split completo \(S=K_P\vee I_O\)
difiere de \(G\) en exactamente \(m+A\) aristas. Por tanto, su grado
máximo de modificación satisface

\[
h=\Delta(G\triangle S)\le m+A\le20d\le n/100.
\tag{8}
\]

Para controlar el tamaño de la raíz, definamos \(Q(n)=(2n+1)^2/24\).
Completando el cuadrado obtenemos

\[
Q(n)-B(n,p)=\frac{(6p-2n-1)^2}{24},
\qquad 0<Q(n)-M(n)<25/24.
\]

La identidad cuadrática mide la distancia del tamaño de la raíz al
centro de la parábola split. Por (7), \(M(n)-B(n,p)\le d\); sumar el
error de parte entera \(Q(n)-M(n)\) y multiplicar por 24 demuestra (4).
La cota de ese error se deduce de \(Q(n)=n(n+1)/6+1/24\) y de la
propiedad definitoria de la parte entera. Como \(d\le n/2000\) y
\(n\ge100\),

\[
24d+25\le0.012n+25\le n^2/100.
\]

Tomando raíces cuadradas obtenemos \(|6p-2n-1|\le n/10\). Así,
\(1.9n+1\le6p\le2.1n+1\). Para \(n\ge100\), estas desigualdades implican
las siguientes cotas gruesas convenientes:

\[
3n/10\le p\le2n/5.
\tag{9}
\]

Definamos \(q=n-p\). Las ecuaciones (8)--(9) dan

\[
q-p=n-2p\ge n/5\ge18h,
\qquad 4p-q=5p-n\ge n/2\ge36h.
\]

Estas son precisamente

\[
p+18h\le q\le4p-36h,
\]

las hipótesis de la Proposición 2.1 de VI con su parámetro
\(t_{\rm par}=4\). En efecto, sus desigualdades son
\(p+(4t_{\rm par}+2)h\le q\) y
\(q+t_{\rm par}(t_{\rm par}+5)h\le t_{\rm par}p\). La positividad
requerida del núcleo y la condición \(q\ge p\) se deducen de (9) y de la
primera desigualdad. La raíz no tiene aristas internas ausentes. Por
tanto, la fórmula local es \(pq-\binom p2-A-2m\), que es (1). Restándola
de \(M(n)\) obtenemos (2). Los tres sumandos de (2) son no negativos, de
modo que \(A+m\le A+2m\le d\), lo que demuestra (3). La cota inicial
\(20d\) solo era un presupuesto de entrada; la igualdad exacta la mejora
a \(d\).

Falta comprobar que la raíz utilizada es el conjunto especificado en el
teorema. Un vértice del clique tiene grado \(n-1\) en la referencia
split completa y pierde a lo sumo \(h\) aristas incidentes. Un vértice
exterior tiene grado de referencia \(p\) y gana a lo sumo \(h\). Por
tanto, para \(a\in P\) y \(x\in O\), las cotas de grado de modificación
dan

\[
d_G(a)\ge n-1-h\ge0.99n-1>n/2,
\qquad d_G(x)\le p+h\le0.41n<n/2.
\]

Así, esta misma raíz es exactamente el conjunto definido por grados en
el enunciado. Esto demuestra el teorema.

\subsection{5. Relación con la serie}\label{relaciuxf3n-con-la-serie}

La existencia de la raíz clique y la exclusión del régimen lejano
provienen del Paper IV. La fórmula exacta en el entorno resultante
proviene del Paper VI. La cota inicial \(A+m\le20d\) paga la entrada en
ese entorno; la igualdad exacta la mejora después a \(A+m\le d\). La
desigualdad abstracta de transferencia del Paper V no proporciona la
raíz estructural. Su aplicación a la cota eventual de Erdős 81 se
desarrolla en la nota IV--V separada.

Cuando \(d=0\), la ecuación (2) fuerza \(A=m=0\) y \(B(n,p)=M(n)\). Por
tanto, el grafo es split completo, con un tamaño de clique maximizante.
Esa clasificación ya está establecida en IV {[}1, Corolario 6.2{]} y en
Okechukwu {[}4, Corolario 1.2{]}. Esta nota registra una identidad
exacta del déficit y el reconocimiento por grados en toda la ventana
lineal enunciada.

Okechukwu {[}4, Teorema 1.3{]} demuestra estabilidad cualitativa para
grafos con defecto simplicial enraizado sublineal; en la Sección 7
plantea cuantificar la relación entre el déficit de partición y la
distancia de edición. El Paper IV {[}1, Teorema 6.1{]} ya da una cota
lineal de edición a un grafo split completo asociado y distingue ese
grafo de uno cuyo núcleo tiene tamaño maximizante. Esta aplicación
combina IV con VI para obtener la cota (3), de coeficiente uno, y la
identidad exacta (2) en la ventana menor \(d\le n/2000\). La distancia
de edición se mide respecto del grafo \(K_P\vee I_O\) identificado por
grados, cuyo tamaño de núcleo no tiene por qué maximizar \(B(n,p)\). Por
tanto, no afirma una tasa cuantitativa óptima para toda la clase
considerada en {[}4{]}. El problema original de partición en cliques de
grafos cordales se debe a Erdős, Ordman y Zalcstein {[}5{]}.

\subsection{Formalización y
reproducibilidad}\label{formalizaciuxf3n-y-reproducibilidad}

El proyecto separado en Lean 4.35 demuestra el Teorema 1.1 a partir de
los insumos demostrados de IV y VI. El companion formal identifica su
paquete final congelado de fuentes, el resultado público y los
adaptadores físicos, el build de la raíz canónica y las auditorías de
axiomas, y los comandos de reproducción con versiones fijadas. Su cadena
seleccionada de dependencias excluye el teorema final de cota global de
IV, conservando sus constructores estructurales y de redondeo. Los
registros de fuentes y auditoría distinguen los ocho módulos de este
proyecto de los cortes exactos de fuentes previas.
\href{https://github.com/jtraverso/erdos-81-chordal-clique-partitions/tree/main/series_notes/NOTE_II_CHORDAL_STABILITY/05_formalization}{Fuentes
y companion formal}.

\subsection{Agradecimientos}\label{agradecimientos}

El autor agradece profundamente a su esposa María Paz y a sus hijos
Lucas, Juan Cristóbal, Francisca, Raimundo y Benjamín su amor, paciencia
y apoyo.

\subsection{Uso de inteligencia artificial y herramientas
computacionales}\label{uso-de-inteligencia-artificial-y-herramientas-computacionales}

En los artículos fuente utilizados aquí, se utilizaron Claude, de
Anthropic, y ChatGPT/Codex, de OpenAI, para explorar argumentos,
comprobarlos y preparar los manuscritos. Aristotle, de Harmonic,
participó en la exploración de caminos, la búsqueda de contraejemplos,
la elaboración de pruebas y su formalización y revisión en Lean; no sólo
tradujo pruebas terminadas.

Para esta nota, ChatGPT/Codex ayudó a ensamblar y comprobar el
argumento, desarrollar las interfaces Lean y preparar el manuscrito y
los materiales de auditoría. Aristotle aportó una propuesta separada de
formalización.

\enlargethispage{2\baselineskip}Certo y Jacobian son las herramientas computacionales usadas en el
programa de investigación subyacente para explorar instancias, comprobar
identidades y evaluar candidatos, según se declara en el Paper IV. Su
uso se declara por separado de las dependencias lógicas de la prueba
formal. Las pruebas finitas de regresión realizadas específicamente para
esta nota utilizan los scripts de aritmética exacta identificados en la
auditoría adjunta. Los cálculos se revisaron según su alcance; una
búsqueda finita no reemplaza una prueba universal.

\Needspace{3\baselineskip}El autor conserva la responsabilidad por los argumentos, las citas, el
código y la presentación. Ningún sistema de IA figura como autor.

\subsection{Referencias}\small\label{referencias}

{[}1{]} J. P. Traverso Gianini, \emph{Clique partitions with rooted
simplicial defect: quantitative stability and sharp eventual bounds},
Paper IV, preprint v1.23, 1 de octubre de 2026.
\href{https://github.com/jtraverso/erdos-81-chordal-clique-partitions/blob/054593443bdece587dda9686a411bb18b0ca34c5/preprints/PAPER_IV/01_manuscript/PAPER_IV_preprint_v1.23_en.md}{Manuscrito
y materiales}. Depósito:
\href{https://doi.org/10.5281/zenodo.23089131}{doi:10.5281/zenodo.23089131}.

{[}2{]} J. P. Traverso Gianini, \emph{Optimal clique partitions under
local perturbations}, Paper VI, Preprint 1.0. La demostración utiliza la
interfaz paramétrica de fuentes congeladas identificada en el companion,
18 de octubre de 2026.
\href{https://github.com/jtraverso/erdos-81-chordal-clique-partitions/tree/main/preprints/PAPER_VI}{Manuscrito
y materiales}.

{[}3{]} J. P. Traverso Gianini, \emph{From realizability certificates to
the Erdős 81 bound}, nota de investigación separada complementaria de
los Papers IV y V, Preprint 1.0, 18 de octubre de 2026.
\href{https://github.com/jtraverso/erdos-81-chordal-clique-partitions/tree/main/series_notes/NOTE_I_ERDOS81_TRANSFER}{Manuscrito
y materiales}.

{[}4{]} O. Okechukwu, \emph{Clique partitions and bounded simplicial
defect}, arXiv:2609.20871v1 (2026).
\url{https://arxiv.org/abs/2609.20871v1}.

{[}5{]} P. Erdős, E. T. Ordman, y Y. Zalcstein, ``Clique partitions of
chordal graphs,'' \emph{Combinatorics, Probability and Computing}
\textbf{2} (1993), 409--415.
\url{https://doi.org/10.1017/S0963548300000808}.

\end{document}
